\chapter{Fonctionnalités optionnelles}

Ces fonctionnalités dépendent d'une source de données ou d'un capteur
supplémentaire. Elles seront traitées une fois la version de base terminée,
dans l'ordre de ce chapitre.

\section{Suggestions selon la météo et tirage au sort}

\begin{itemize}
  \item \textbf{Carte du jour selon la météo.} Le bandeau de l'accueil
        affiche la date et le temps local (écran 1), et le tirage privilégie
        les plats froids quand il fait chaud. Il faut une API météo sans clé,
        comme Open-Meteo, et la localisation approximative.
  \item \textbf{Le verre du hasard.} Un cocktail tiré au sort sur l'accueil
        (écran 2), renouvelé en secouant le téléphone (accéléromètre).
\end{itemize}

\section{Informations nutritionnelles et scan}

\begin{itemize}
  \item \textbf{Nutri-Score estimé} d'un plat (écran 3b). TheMealDB ne
        fournit aucune valeur nutritionnelle et exprime les quantités en
        texte libre. Le calcul suppose une table de composition des
        aliments et une conversion des mesures en grammes. Le résultat
        reste une estimation, présentée comme telle.
  \item \textbf{Scan d'un code-barres} en magasin (écran 18) : CameraX et
        ML Kit pour la lecture, Open Food Facts pour la fiche produit.
        Un produit reconnu peut être coché directement dans la liste de
        courses.
\end{itemize}

\section{Variantes sans alcool}

Le sélecteur « Version sans alcool » d'une fiche cocktail (écran 13). Aucune
API ne relie un cocktail à son équivalent sans alcool : il faut une table de
correspondances rédigée par l'équipe.

\section{Planification des repas}

Associer des menus aux jours de la semaine et générer la liste de courses de
la semaine entière.

\section{Recherche par ingrédients disponibles}

L'utilisateur indique ce qu'il a dans ses placards et obtient les recettes
réalisables. L'API gratuite ne filtre que sur un ingrédient à la fois :
l'application croise elle-même les résultats.

\section{Estimation du coût}

Aucune des deux API ne donne de prix. L'estimation reposerait sur une table
de prix moyens par ingrédient, forcément indicative.

\section{Traduction des recettes}

Les contenus des API sont en anglais. Une traduction sur l'appareil (ML Kit
Translation) rendrait les fiches en français. Dans la version de base, seule
l'interface est en français.
